%% Reference/navigation apparatus; the registry owns every rule display.
\chapter{Nonclassical and Specialized Calculus Sheets}
\label{sec:p14-nonclassical-and-specialized-calculi}

A calculus is selected by its judgments and its permitted uses of premises,
not just by the name of a logic.  This chapter links the registered rule
surfaces to their language, context, side conditions, semantics, and exact
adequacy claims.  The links supplement the full treatments; they do not
identify different presentations or replace their proofs.  In particular,
soundness, completeness, cut elimination, proof normalization, and termination
of a search procedure are different properties.  The presence of a Cut row
does not establish its admissibility in the cut-free subsystem, and an
effective normal form for formulas is not a normalization theorem for every
proof \cite{Avron1996,Negri2005,DunnRestall2002,Pitts1995Draft}.

In the modal Hilbert rows, \(A,B,C\) are formulas and \(\Gamma\) is a set of
local formula premises.  In the temporal rows, state formulas and path
formulas have their explicitly separate grammars.  In the resource rows,
\(\Gamma,\Delta\) denote the displayed multiset or ordered-sequence
contexts.  In the contextual first-order rows, \(\Gamma\) instead lists
distinct sorted variables and \(\Phi\) lists formula hypotheses.  The same
written context letter does not license transfer between these judgment
types.  Each sheet retains its own source notation.

\section{Modal, temporal, epistemic, dynamic, and hybrid systems}
\label{sec:p14-nonclassical-modal-temporal}

\subsection{Normal Hilbert systems and local consequence}

The common normal base is Sheet~\ref{sheet:hilbert-k}: classical propositional
tautology instances, the distribution schema K, modus ponens, and
\emph{theorem-only} necessitation.  A derivation from \(\Gamma\) may use a
formula \(\Box A\) obtained by necessitation only when \(A\) also has a
separate derivation with no premises.  It may not box an arbitrary line whose
derivation depends on \(\Gamma\).  Thus local consequence assumes the
premises at the evaluation world, not at every accessible world.  The
countermodel and derivation distinction are explicit in
Lemma~\ref{thm:p08-mce-separations} and
Theorem~\ref{thm:p08-modal-canonical-completeness}.

\begingroup\small
\begin{longtable}{@{}p{.16\linewidth}p{.29\linewidth}p{.47\linewidth}@{}}
\caption{Normal modal sheet identities and their controlling adequacy routes.}
\label{tab:p14-nonclassical-normal-modal}\\
\toprule
System and sheet & Extension and exact frame regime & Adequacy and proof route \\
\midrule\endfirsthead
\caption[]{Normal modal sheet identities and their controlling adequacy routes (continued)}\\
\toprule System and sheet & Extension and exact frame regime & Adequacy and proof route \\
\midrule\endhead
\midrule\multicolumn{3}{r}{\emph{continued on next page}}\\
\endfoot
\endlastfoot
\(K\), \ref{sheet:hilbert-k} & Common normal base; arbitrary nonempty
Kripke frames. & Strong local completeness for arbitrary premise sets is
proved in Theorem~\ref{thm:p08-modal-canonical-completeness},
Annex~\ref{annex:p8.modal.1}. The separate inherited four-system statement
retains its sketch boundary in Theorem~\ref{thm:p08-modal-source-sketch}. \\
\addlinespace
\(KD\), \ref{sheet:hilbert-d} & K plus D; serial frames. & Strong local
completeness over the declared countable language is
Theorem~\ref{thm:p08-kd-strong-completeness},
Annex~\ref{annex:p8.mce.7}. \\
\addlinespace
\(KT\), \ref{sheet:hilbert-t} & K plus T; reflexive frames. & The same
complete four-system canonical proof applies to this exact sheet, not to
every reflexive modal extension. \\
\addlinespace
\(KTB\), \ref{sheet:hilbert-b} & KT plus B; reflexive symmetric frames. &
Theorem~\ref{thm:p08-ktb-strong-completeness},
Annex~\ref{annex:p8.mce.8}. The sheet is not the symmetry-only system KB;
the source-specific alias B is explained in
Warning~\ref{def:p08-modal-b-convention}. \\
\addlinespace
\(K4\), \ref{sheet:hilbert-k4} & K plus 4; transitive frames. &
Theorem~\ref{thm:p08-k4-strong-completeness},
Annex~\ref{annex:p8.mce.9}, proves strong local completeness. \\
\addlinespace
\(S4\), \ref{sheet:hilbert-s4} & KT plus 4; reflexive transitive frames. &
Theorem~\ref{thm:p08-modal-canonical-completeness} supplies the complete
canonical proof. The source name KT4 is preserved. \\
\addlinespace
\(KD45\), \ref{sheet:hilbert-kd45} & KD plus 4 and 5; serial transitive
Euclidean frames. & Strong local completeness is source-cited in
Entry~\ref{entry:p08-modal-kd45}, \cite{ChagrovZakharyaschev1997}. \\
\addlinespace
\(S4.2\), \ref{sheet:hilbert-s42} & S4 plus .2; directed S4 frames. &
Entry~\ref{entry:p08-modal-s42} cites strong local completeness for that
directed frame class, \cite{ChagrovZakharyaschev1997}. \\
\addlinespace
\(S4.3\), \ref{sheet:hilbert-s43} & S4 plus .3; weakly connected S4
frames. & Entry~\ref{entry:p08-modal-s43} cites the corresponding strong
local result. Directedness and weak connectedness are the distinct clauses
in Definition~\ref{def:p08-modal-frame-properties}. \\
\addlinespace
\(S5\), \ref{sheet:hilbert-s5} & KT plus 5; equivalence frames. &
Theorem~\ref{thm:p08-modal-canonical-completeness} proves strong local
completeness. The sheet follows the source's T-plus-5 presentation; 4 is
derivable, not an extra primitive row of this sheet. \\
\addlinespace
\(\mathsf{GL}\), \ref{sheet:hilbert-gl} & K plus L\"ob; finite transitive
irreflexive frames for theoremhood completeness. &
Entry~\ref{entry:p08-modal-gl} cites the relational result; the arithmetical
realization and Solovay theorem are
Definition~\ref{def:p05-arithmetical-realization} and
Theorem~\ref{thm:p05-solovay-completeness}. Neither is the ordinary
arbitrary-frame canonical argument. \\
\addlinespace
\(\mathsf{Grz}\), \ref{sheet:hilbert-grz} & S4 plus the Grzegorczyk schema;
finite partial-order frames for theoremhood completeness. &
Entry~\ref{entry:p08-modal-grz} gives the cited result and preserves its
finite-frame theoremhood boundary, \cite{ChagrovZakharyaschev1997}. \\
\bottomrule
\end{longtable}\endgroup

These Hilbert sheets do not formulate a sequent Cut rule or a reduction
relation on proofs.  Their completeness results therefore are not recorded
as proof-normalization claims.  Frame correspondence, a finite-model
property, and completeness also have separate hypotheses: the exact
correspondences are Table~\ref{tab:p08-modal-correspondence}, and the
filtration construction is Section~\ref{sec:p08-modal-filtration}.

\showcalculus{hilbert-k}
\showcalculus{hilbert-d}
\showcalculus{hilbert-t}
\showcalculus{hilbert-b}
\showcalculus{hilbert-k4}
\showcalculus{hilbert-s4}
\showcalculus{hilbert-kd45}
\showcalculus{hilbert-s42}
\showcalculus{hilbert-s43}
\showcalculus{hilbert-s5}
\showcalculus{hilbert-gl}
\showcalculus{hilbert-grz}

\subsection{Non-normal systems and alternative modal derivation objects}

For E, EM, EC, EN, and EMC, the common base is
Sheet~\ref{sheet:hilbert-e}, with theorem-only replacement of provable
equivalents rather than the normal K/Nec package.  M adds supplementation,
C adds binary intersection closure, and N adds the unit neighborhood.  The
five existing sheets are respectively
\ref{sheet:hilbert-e}, \ref{sheet:hilbert-m}, \ref{sheet:hilbert-c},
\ref{sheet:hilbert-n}, and \ref{sheet:hilbert-r}.  The last is the regular
modal logic EMC, not the relevant implication calculus \(\mathbf R\).
The source's primitive-possibility language and the sheets' defined-diamond
language are connected by the explicit translation in
Definition~\ref{def:p08-cn-language} and
Theorem~\ref{thm:p08-cn-language}; they are not silently identical grammars.

The source rule/axiom correspondence proves that Chellas's system M is EM,
R is EMC, and the normal base is EMCN
(Theorem~\ref{thm:p08-cn-rules}).  EN alone is not K, and EMC does not already
contain N: the finite witnesses are
Theorem~\ref{thm:p08-cn-separation}.  For all eight M/C/N combinations, the
complete neighborhood construction and strong local-consequence extension
are Theorem~\ref{thm:p08-cn-completeness} and
Annex~\ref{annex:p8.cn.completeness}.  The model closure conditions concern
sets of worlds, while RE concerns an independently proved formula
equivalence.  Agreement at one world cannot be substituted under a modality.
These distinctions are needed even when the written Boolean rules coincide
\cite[chaps.~7--9]{Chellas1980}.

\begingroup\small
\begin{longtable}{@{}p{.20\linewidth}p{.34\linewidth}p{.38\linewidth}@{}}
\caption{Modal architectures: contexts, restrictions, and scope of cut claims.}
\label{tab:p14-nonclassical-modal-architectures}\\
\toprule
Exact sheet & Derivation object and side conditions & Worked and metatheoretic targets \\
\midrule\endfirsthead
\caption[]{Modal architectures: contexts, restrictions, and scope of cut claims (continued)}\\
\toprule Exact sheet & Derivation object and side conditions & Worked and metatheoretic targets \\
\midrule\endhead
\midrule\multicolumn{3}{r}{\emph{continued on next page}}\\
\endfoot
\endlastfoot
Avron GS5, \ref{sheet:gs5-avron} & Finite-multiset hypersequents, with
classical LK rules in one component and external context retained.
Box-R has an entirely box-prefixed antecedent; internal and external
weakening and external contraction are explicit. & Axiom 5 is derived in
Example~\ref{def:p07-gs5-derivation}. Adequacy and cut elimination are
source-cited in Theorem~\ref{thm:p07-gs5-adequacy}, \cite{Avron1996}. \\
\addlinespace
Negri labelled K, \ref{sheet:labelled-k-negri} & Finite-multiset sequents
of labelled formulas and relational atoms. Labels are metalinguistic world
variables; Box-R requires its eigenlabel fresh for the conclusion.
The K sheet has no added reflexivity or transitivity frame rule. &
Example~\ref{def:p07-labelled-k-derivation} derives K.
Theorem~\ref{thm:p07-labelled-k-adequacy} cites soundness, completeness,
and cut admissibility, \cite{Negri2005}. \\
\addlinespace
Wansing DK fragment, \ref{sheet:display-k-wansing} & Display sequents
in a structural language with unit, structural negation, product, and modal
operator. Display postulates are reversible; a formula rule requires the
principal formula displayed. &
Example~\ref{def:p07-display-k-derivation} uses precisely this fragment.
Theorem~\ref{thm:p07-display-k-adequacy} cites adequacy and cut for
Wansing's \emph{full} calculus, not for the printed fragment in isolation,
\cite{Wansing1994}. \\
\addlinespace
Stouppa \(\mathsf{KSg\mbox{-}S5}\), \ref{sheet:deep-s5-stouppa} &
Negation-normal structures modulo the stated congruences. Every rule applies
inside a positive structure context; proofs start at \(\mathsf{tt}\).
Brackets represent disjunction and parentheses conjunction, not a sequent
antecedent. & Example~\ref{def:p07-deep-s5-derivation} derives B by
non-root rule applications. The cited S5 adequacy and cut-admissibility
result is Theorem~\ref{thm:p07-deep-s5-adequacy}, \cite{Stouppa2007}. \\
\bottomrule
\end{longtable}\endgroup

\showcalculus{hilbert-e}
\showcalculus{hilbert-m}
\showcalculus{hilbert-c}
\showcalculus{hilbert-n}
\showcalculus{hilbert-r}
\showcalculus{gs5-avron}
\showcalculus{labelled-k-negri}
\showcalculus{display-k-wansing}
\showcalculus{deep-s5-stouppa}

\subsection{Time, paths, and fixed-point judgments}

The temporal sheets retain both their temporal convention and their judgment
type.  LTL uses discrete infinite words and strong until with a witness at
the current or a later position.  The strict-future source convention is
connected to it only by the translations proved in
Theorem~\ref{thm:p08-ltl-strict-bridge}.  CTL and CTL* use total transition
structures, so their paths are infinite.  Removing totality or changing to
finite traces is a change of semantics, not a harmless implementation choice.
Their state/path grammars are Definitions~\ref{def:p08-ctl-semantics}
and~\ref{def:p08-ctl-star-semantics}.

\begingroup\small
\begin{longtable}{@{}p{.16\linewidth}p{.36\linewidth}p{.40\linewidth}@{}}
\caption{Temporal and fixed-point sheets: rule-use restrictions and completeness.}
\label{tab:p14-nonclassical-temporal}\\
\toprule
Sheet & Rules and exact judgment & Proof and semantic scope \\
\midrule\endfirsthead
\caption[]{Temporal and fixed-point sheets: rule-use restrictions and completeness (continued)}\\
\toprule Sheet & Rules and exact judgment & Proof and semantic scope \\
\midrule\endhead
\midrule\multicolumn{3}{r}{\emph{continued on next page}}\\
\endfoot
\endlastfoot
LTL, \ref{sheet:hilbert-ltl} & No-premise Hilbert theoremhood.
XNec and UInd require no-premise derivations of their displayed premises;
U-unfolding alone does not settle eventualities. & The full finite-type,
eventuality, realization and truth-lemma development is
Section~\ref{sec:p08-linear-temporal-completeness} and
Annex~\ref{annex:p8.ltl.source}. The arbitrary-premise failure is proved in
Theorem~\ref{thm:p08-ltl-noncompactness}; weak validity completeness is
not strong local completeness. \\
\addlinespace
CTL, \ref{sheet:hilbert-ctl} & State-formula theoremhood.
EX-Mon and both until-induction rules require theorem premises; the
Total axiom belongs to the selected transition regime. &
Entry~\ref{entry:p08-ctl} cites weak completeness for total structures,
\cite[\S6]{EmersonHalpern1985}. Correctness of finite-state CTL labelling
has a separate full proof at Theorem~\ref{thm:p09-ctl-labelling},
Annex~\ref{annex:p9.ctl.1}; that is an algorithm theorem, not the Hilbert
completeness proof. \\
\addlinespace
Reynolds CTL*, \ref{sheet:hilbert-ctl-star} & State and path formulas
remain typed. TG, PG and AA are theorem-only. C14 applies only to an atom;
there is no unrestricted substitution inference. AA retains its finite
auxiliary-atom, partition, and functional-expansion conditions. &
Entry~\ref{entry:p08-ctl-star} cites weak completeness for standard
fullpath validity, including LC and AA, \cite{Reynolds2001}. The bundled
base without those additions is not the same standard-fullpath calculus;
arbitrary-premise strong completeness is not claimed. \\
\addlinespace
Kozen \(L_\mu\), \ref{sheet:hilbert-mu} & Boolean equational judgments,
with inequality abbreviating join equality and theoremhood meaning equality
to top. Bound fixed-point variables occur positively; all substitutions
avoid capture. Park requires its displayed fresh variable absent from the
body. & The exact syntax and powerset semantics are
Definition~\ref{def:p08-modal-mu-semantics}. Weak completeness of this
equational presentation is Theorem~\ref{thm:p09-kozen-full-completeness},
source-cited to \cite{Walukiewicz2000}; it is not a normalization theorem
or arbitrary-premise standard strong completeness. \\
\bottomrule
\end{longtable}\endgroup

Least- and greatest-fixed-point semantics act on sets of states; program
constants index relations, while fixed-point variables range over state
sets.  This typing prevents identifying Park induction with UInd merely
because both express an induction principle.  The temporal and fixed-point
comparisons are Table~\ref{tab:p08-temporal-comparison}; the actual
finite-model and decision constructions retain their own proof targets.

\showcalculus{hilbert-ltl}
\showcalculus{hilbert-ctl}
\showcalculus{hilbert-ctl-star}
\showcalculus{hilbert-mu}

\subsection{Knowledge, belief, announcements, programs, and named worlds}

The multi-agent sheets are componentwise: there is one modality per agent
and no cross-agent interaction axiom.  Knowledge
(Sheet~\ref{sheet:hilbert-s5n}) uses equivalence relations and includes
factivity; belief (Sheet~\ref{sheet:hilbert-kd45n}) uses serial transitive
Euclidean relations and does not infer factivity.  Necessitation is
theorem-only in both premise derivability conventions.  The epistemic
adequacy statement remains source-sketched at
Entry~\ref{entry:p08-s5n}, Annex~\ref{annex:p8.s5n.source}; the doxastic
local strong-completeness result is cited at
Entry~\ref{entry:p08-kd45n} to \cite{FaginHalpernMosesVardi1995}.

Sheet~\ref{sheet:hilbert-common-knowledge} adds common knowledge for a fixed
nonempty finite group, retaining a theorem-only induction rule.
Everyone-knows is a finite conjunction, distributed knowledge uses the
intersection of the agent relations, and common knowledge uses the
reflexive--transitive closure of their union.  Their precise separation is
Definition~\ref{def:p08-distributed-knowledge},
Example~\ref{def:p08-everyone-not-common}, and
Table~\ref{tab:p08-epistemic-systems}.  The weak completeness claim and the
failure of arbitrary-premise strong completeness are kept separate in
Entry~\ref{entry:p08-common-knowledge} and
Theorem~\ref{thm:p08-common-knowledge-noncompact}; the latter needs at least
two agents.  A group consisting of one S5 agent does not witness that failure.

Public-announcement reduction on Sheet~\ref{sheet:hilbert-pal} preserves its
exact S5 base.  The atom row is atomic-only, and replacement under an
announcement is through a theorem equivalence.  An announcement restricts
the model to its precondition worlds; its box is vacuously true when the
precondition is false.  The cited reduction/completeness target is
Theorem~\ref{thm:p08-pal-reduction}, \cite{Plaza1989}.  Reduction to static
epistemic formulas is not a general event-model reduction or a proof
normalization theorem.  The true-but-unsuccessful announcement in
Example~\ref{def:p08-unsuccessful-announcement} shows why truth before update
does not imply truth after it.

PDL theoremhood uses Sheet~\ref{sheet:hilbert-pdl}.  Programs and formulas
are generated together: composition, choice, star, and tests are program
constructors, and necessitation by a program is theorem-only.  The
test-containing presentation is attributed to
\cite{HarelKozenTiuryn2000}, while \cite{KozenParikh1981} supplies the distinct
test-free core.  The completeness target is
Theorem~\ref{thm:p09-pdl-adequacy}; Fischer--Ladner closure and decision
claims have the separate target Theorem~\ref{thm:p09-pdl-decision}.  Kleene
star denotes reflexive--transitive closure of a relation, not unrestricted
program recursion.

Basic hybrid \(K(@)\) uses Sheet~\ref{sheet:tableaux-hybrid-k}, not the
metalinguistic label calculus for K above.  Nominals belong to the
object language and each names exactly one world.  A satisfiability tableau
starts with a positive satisfaction formula at a fresh nominal; a validity
refutation starts with its negative counterpart.  Fresh witness nominals,
one-time expansion, nominal reflexivity, and uniform nominal replacement
are part of the registered procedure.  Nominal identity is not ordinary
first-order object equality.  The cited adequacy target is
Theorem~\ref{thm:p08-hybrid-tableau}, \cite{ArecesTenCate2007}; binders and
the global modality are excluded from this exact language.

\showcalculus{hilbert-s5n}
\showcalculus{hilbert-kd45n}
\showcalculus{hilbert-common-knowledge}
\showcalculus{hilbert-pal}
\showcalculus{hilbert-pdl}
\showcalculus{tableaux-hybrid-k}

\section{Conditional, deontic, free, and inquisitive systems}
\label{sec:p14-nonclassical-conditional-inquisitive}

\subsection{Conditional theorem rules and obligation readings}

Chellas CK on Sheet~\ref{sheet:hilbert-conditional-ck} is a theoremhood
calculus over a classical Boolean base.  RCEA replaces antecedents only
through a theorem equivalence, and RCK lifts a theorem finite-conjunction
implication to conditionals sharing one antecedent.  Its zero-conjunct case
is included.  Neither rule permits conditionalizing an open premise merely
because it is present in a local derivation.  CK's set-valued selection
function need not select antecedent worlds; adding that condition validates
ID and changes the system.  The exact semantics and cited completeness
are Definition~\ref{def:p08-ck-semantics} and
Theorem~\ref{thm:p08-ck-completeness}, \cite{Chellas1975}.

Lewis VC, also called C1 in the cited conditional nomenclature, extends CK
on Sheet~\ref{sheet:hilbert-counterfactual-vc} by ID, CA, CSO, CV,
conditional MP, and CS.  The displayed corrected Nute-style presentation
uses CA rather than the non-equivalent MOD schema.  Its controlling
equivalence/source roles are retained on the sheet
\cite{Lewis1971,Lewis1973,Wen2019}.  Centered sphere semantics is
Definition~\ref{def:p08-lewis-conditional}; the cited completeness result is
Theorem~\ref{thm:p08-vc-completeness}.  Centering does not add the limit
assumption, and this VC presentation permits several equally close
antecedent worlds.  It is not Stalnaker's unique-selection system, nor is
conditional C1 da Costa's paraconsistent \(C_1\).

Standard deontic logic is the exact monadic KD system of
Definition~\ref{def:p08-sdl}, with the box read as obligation and its dual
as permission.  Its calculus is Sheet~\ref{sheet:hilbert-d}; seriality,
not reflexivity, controls D.  The Ross inference is a valid consequence
of normality in that system, not an additional primitive inference rule.
The contrary-to-duty and scope diagnoses are
Table~\ref{tab:p08-deontic-systems}.  The dyadic best-world semantics in
Definition~\ref{def:p08-dyadic-obligation} instead chooses the best accessible
condition-worlds under a smooth total preorder and requires that set to be
nonempty.  An impossible condition therefore makes this selected dyadic
obligation false.  It is not monadic KD with a material implication inserted
inside the box, and no KD completeness theorem is transferred to it
\cite{Hansson1969,McNamara2006}.

\showcalculus{hilbert-conditional-ck}
\showcalculus{hilbert-counterfactual-vc}

\subsection{Free-logic guards are variant-indexed}

Free-logic rules must fix ordinary term denotation, failed descriptions,
atomic truth, equality, the quantifier domain, and designation before an
existence guard can be interpreted.  The three selected variants in
Table~\ref{tab:p08-free-logic-variants} do not form one undifferentiated
classical first-order calculus.  Their semantic consequence and guarded
quantifier surfaces are Definition~\ref{def:p08-free-logic-consequence} and
Table~\ref{tab:p08-free-logic-quantifier-rules}.

\begingroup\small
\begin{longtable}{@{}p{.15\linewidth}p{.39\linewidth}p{.38\linewidth}@{}}
\caption{Free-logic quantifier interface: why the same guard has different standing.}
\label{tab:p14-nonclassical-free-guards}\\
\toprule
Variant & Term, atom, equality, and quantifier policy & Permitted use of the guarded semantic rules \\
\midrule\endfirsthead
\caption[]{Free-logic quantifier interface: why the same guard has different standing (continued)}\\
\toprule Variant & Term, atom, equality, and quantifier policy & Permitted use of the guarded semantic rules \\
\midrule\endhead
\midrule\multicolumn{3}{r}{\emph{continued on next page}}\\
\endfoot
\endlastfoot
Negative & Ordinary terms denote in the outer domain; predicates are
strict on the existing inner domain, and identity requires both values to
exist. A failed primitive description makes its containing atom false.
Quantifiers range over the inner domain. & Universal instantiation and
existential generalization have the explicit existence premise from the
linked table. Outer denotation alone is insufficient. The negative clause
must not be replaced by a positive predication clause. \\
\addlinespace
Positive & Ordinary terms have outer denotation; predicate extensions may
include nonexisting objects, and identity is outer identity. Quantifiers
still range over the inner domain. & A true atomic predication does not
itself supply existence. The same written existence guard remains needed
for the displayed quantifier consequences. \\
\addlinespace
Neutral & Terms may be partial; an undefined argument gives its atom or
identity the strong-Kleene gap value. Only true is designated; quantifiers
use strong-Kleene infima and suprema over the inner domain. & Both
premises of a displayed guarded consequence must be designated. A gap or
nonexisting value of the guard licenses no step. Classical two-valued
reasoning cannot be imported solely from the guard's spelling. \\
\bottomrule
\end{longtable}\endgroup

Russellian contextual elimination of a description and a primitive
description term are separate presentations, as the full treatment explains
\cite{Russell1905,Lambert1967,Bencivenga1986,Lehmann2002}.  The guarded
table asserts semantic consequence under these fixed variants; it does not
assert interderivability, cut elimination, or completeness of a newly named
family-wide proof calculus.

For the selected negative convention, the separately derived certificate
calculus is \cref{sec:p08-negative-free-logic-calculus} and
Sheet~\ref{sheet:negative-free-logic}. Its defining premise is an actual
finite proof in the registered classical target language, not a semantic
validity oracle. \Cref{thm:p08-negative-model-bridge,thm:p08-negative-adequacy}
prove encoding/decoding and strong adequacy. The native guarded rows are
proved interfaces of that certificate presentation; they are not asserted
to form an independently complete native calculus on their own.
Free parameters are outer-valued unless the named existing-parameter
context requires otherwise. Strictness applies to complete atomic
argument values: a total ordinary function may map a nonexisting input to
an existing output. The partial-description result uses the explicit strict
undefined-input lift, and neither the positive nor neutral variant is
silently covered by this negative proof.

\showcalculus{negative-free-logic}

\subsection{Inquisitive deduction and assumption discharge}

Sheet~\ref{sheet:p08-inquisitive} presents natural deduction for
\(\mathsf{InqB}\), with arbitrary formulas distinguished from syntactically
classical formulas.  Its inquisitive disjunction selects a disjunct supported
by the whole information state; it is not the derived classical
disjunction.  Implication introduction and inquisitive-disjunction
elimination discharge selected assumption occurrences, possibly none, while
the remaining assumptions remain undischarged.  Split requires a classical
antecedent, and double-negation elimination requires a classical formula.
The rules operate from open assumptions, unlike the theorem-only modal and
conditional rules above.

The support clauses and proof definition are
Definitions~\ref{def:p08-inq-support} and~\ref{def:p08-inq-proofs}.
Full soundness is Theorem~\ref{thm:p08-inq-soundness},
Annex~\ref{annex:p8.inq.soundness}; full completeness for arbitrary premise
sets in the fixed countable language is
Theorem~\ref{thm:p08-inq-completeness},
Annex~\ref{annex:p8.inq.completeness}.  The resolution algorithm for a proof
is separately constructed in Theorem~\ref{thm:p08-inq-proof-resolution},
Annex~\ref{annex:p8.inq.proof-resolution}.  These results concern classical
atomic semantics, not arbitrary inquisitive propositions assigned to atoms
\cite[\S\S3.5--3.8]{Ciardelli2023InquisitivePropositional}.

Substitution is legitimate only under its stated classical/flat restriction
(Theorem~\ref{thm:p08-inq-substitution}).  Formula-resolution normal form and
proof-to-answer extraction are the supplied effective results; they do not
by themselves assert a reduction normalization theorem for all proof trees.
Unrestricted classical-disjunction elimination with an inquisitive
conclusion, unrestricted double-negation elimination for questions, and
unrestricted uniform substitution would change the calculus and are not
on this sheet.

\showcalculus{p08-inquisitive}

\section{Many-valued, fuzzy, relevant, and paraconsistent systems}
\label{sec:p14-nonclassical-matrices-relevance}

\subsection{One tableau procedure, separately fixed matrices}

Sheet~\ref{sheet:tableaux-matrix} takes a fixed nonempty finite value set,
a designated subset, and total finite-signature connective operations.
For a finite premise set and one conclusion it creates a separate root for
every assignment of designated premise labels and a nondesignated conclusion
label.  Branches retain their formula occurrences; expansion marks prevent
expanding the same occurrence twice along a branch.  A connective expansion
branches over its complete operation preimage, while a clash uses two
different labels for the \emph{same atom}.  Empty operation preimages close
a branch.  An empty root family is closed, including the degenerate
designation cases allowed by the sheet.

The repaired complete adequacy and saturation proof is in
Section~\ref{sec:many-valued-paraconsistent-02}; its precise specialized
proof targets are recorded in the system entries below.  Changing only
designation can change consequence even when the value set and all operations
are unchanged.

\begingroup\small
\begin{longtable}{@{}p{.16\linewidth}p{.37\linewidth}p{.39\linewidth}@{}}
\caption{Finite-matrix specializations, all using the complete value-labelled sheet.}
\label{tab:p14-nonclassical-matrices}\\
\toprule
Specialization & Fixed language, operations, and designation & Exact adequacy and limitation \\
\midrule\endfirsthead
\caption[]{Finite-matrix specializations, all using the complete value-labelled sheet (continued)}\\
\toprule Specialization & Fixed language, operations, and designation & Exact adequacy and limitation \\
\midrule\endhead
\midrule\multicolumn{3}{r}{\emph{continued on next page}}\\
\endfoot
\endlastfoot
K3 & Constant-free negation, conjunction, disjunction; strong-Kleene
three-value operations, only true designated. &
Entry~\ref{entry:p08-k3}, Annex~\ref{annex:p8.k3.1}, supplies the finite-premise
tableau proof. The absence of tautologies is language-specific, not a
claim about arbitrary truth-constant expansions. \\
\addlinespace
LP & The same operations and fixed language as K3; both and true
designated. & Entry~\ref{entry:p08-lp}, Annex~\ref{annex:p8.lp.1}.
Contradictory premises need not explode. The K3--LP duality is the finite
multiple-conclusion statement in Theorem~\ref{thm:p08-k3-lp-facts},
Annex~\ref{annex:p8.matrix.2}, not an unqualified equality of consequence. \\
\addlinespace
FDE & The four positive/negative support pairs, designating exactly positive
support; the operations are those of
Definition~\ref{def:p08-fde-matrix}. & Entry~\ref{entry:p08-fde},
Annex~\ref{annex:p8.fde.1}. Both and neither are distinct; truth order and
information order are distinct in Definition~\ref{def:p08-fde-orders}. \\
\addlinespace
Finite \(\mathsf L_n\) & The fixed \(n\)-element MV chain, \(n\geq2\),
with implication and negation and only value one designated. &
Lemma~\ref{thm:p08-finite-lukasiewicz-tableau},
Annex~\ref{annex:p8.lukasiewicz.1}, proves terminating finite-premise
adequacy. The number of values and the source chain-index convention are
Definition~\ref{def:p08-lukasiewicz-matrices}. \\
\addlinespace
Bochvar external and fixed Post matrices & Their connective signatures
and operation tables are separately fixed by
Definitions~\ref{def:p08-bochvar-external} and~\ref{def:p08-post-n-valued}. &
Theorem~\ref{thm:p08-bochvar-tableau}, Annex~\ref{annex:p8.bochvar.1},
and Lemma~\ref{thm:p08-post-matrix-decision}, Annex~\ref{annex:p8.post.1},
apply the finite procedure. Infectious weak-Kleene operations alone do not
specify Bochvar's richer external language. \\
\bottomrule
\end{longtable}\endgroup

All these instances use value-preservation consequence, not numeric
averaging or a probabilistic threshold.  The finite exhaustive procedure
cannot be run unchanged over the real interval.  Its termination proof is
an expansion/search result, not a cut-elimination or proof-normalization
result for an unrelated Hilbert calculus \cite{Kleene1952,Priest1979,
OmoriWansing2017,CignoliDOttavianoMundici2000}.

\showcalculus{tableaux-matrix}

\subsection{Infinite-valued and t-norm Hilbert presentations}

Sheet~\ref{sheet:hilbert-lukasiewicz-infinite} has primitive implication and
negation, four axiom schemas, and modus ponens.  Its standard values lie in
the real unit interval; the finite-chain tableau is not its calculus.
Theoremhood and finite-premise value-one completeness are the cited
Theorem~\ref{thm:p08-lukasiewicz-completeness}, \cite{Hajek1998}.
McNaughton's representation result is the separate
Theorem~\ref{thm:p08-mcnaughton}; representing truth functions does not
automatically give arbitrary-premise standard-matrix completeness.

The \hyperref[sheet:hilbert-tnorm-logics]{t-norm collection} has four
separately named Hilbert surfaces.  Lattice conjunction and strong
conjunction are distinct constructors.  The residuum is a truth-function
operation on numeric arguments, not a sequent arrow.  The exact completeness
comparison is Theorem~\ref{thm:p08-tnorm-completeness} and
Table~\ref{tab:p08-tnorm-logics}.

\begingroup\small
\begin{longtable}{@{}p{.16\linewidth}p{.36\linewidth}p{.40\linewidth}@{}}
\caption{Fuzzy Hilbert sheets: extension identity and standard-completeness domain.}
\label{tab:p14-nonclassical-fuzzy}\\
\toprule
Sheet & Exact extension and standard interpretation & Completeness boundary \\
\midrule\endfirsthead
\caption[]{Fuzzy Hilbert sheets: extension identity and standard-completeness domain (continued)}\\
\toprule Sheet & Exact extension and standard interpretation & Completeness boundary \\
\midrule\endhead
\midrule\multicolumn{3}{r}{\emph{continued on next page}}\\
\endfoot
\endlastfoot
MTL, \ref{sheet:hilbert-mtl} & The full displayed MTL axiom menu and MP;
the standard class ranges over all left-continuous t-norms. & Strong
standard completeness for arbitrary premises ranges over that whole class,
not one arbitrarily selected chain, \cite{EstevaGodo2001,JeneiMontagna2002}. \\
\addlinespace
BL, \ref{sheet:hilbert-bl} & MTL plus divisibility; the standard class
ranges over continuous t-norms. & Finite-premise standard completeness;
arbitrary-premise algebraic completeness has a different algebra-class
quantifier, \cite{CignoliEstevaGodoTorrens2000}. \\
\addlinespace
G\"odel G, \ref{sheet:hilbert-godel-fuzzy} & BL plus its contraction
schema; standard strong conjunction is minimum. & Arbitrary-premise
standard completeness for the G\"odel chain is retained separately from
the corresponding algebraic result, \cite{Hajek1998,CintulaHajekNoguera2011}. \\
\addlinespace
Product \(\Pi\), \ref{sheet:hilbert-product-fuzzy} & BL plus both
Product axioms; standard strong conjunction is multiplication. &
Finite-premise standard completeness; no arbitrary-premise standard-chain
claim is imported from MTL or G, \cite{Hajek1998,CintulaHajekNoguera2011}. \\
\bottomrule
\end{longtable}\endgroup

These are source-cited adequacy results, not local cut-elimination proofs.
The algebraic terminology in Section~\ref{sec:p10-many-substructural}
must be read with its stated operations and algebra class; adding a sound
algebra-valued interpretation is not by itself an algebraization theorem.

\showcalculus{hilbert-lukasiewicz-infinite}
\showcalculus{hilbert-tnorm-logics}

\subsection{Relevant theoremhood and paraconsistent mechanisms}

Sheets~\ref{sheet:relevance-e}, \ref{sheet:relevance-r}, and
\ref{sheet:relevance-rm} use the constant-free relevant language, theoremhood
as closure under modus ponens and adjunction, and the fixed Dunn--Restall
schema menu.  R adds Demodaliser to E; RM adds Mingle to R.  Relevant
\(\Box A\) is an eliminable implication abbreviation, not the primitive
normal-modal box.  These sheets do not have a generic antecedent-context
profile that can be filled with classical weakening or contraction rules.

The matched pointed Routley--Meyer classes, hereditary valuations, star,
and validity at the distinguished point are
Definition~\ref{def:p08-pointed-rm-frame}.  The cited adequacy statement is
\emph{formula theoremhood} for those matched classes
(Theorem~\ref{thm:p08-relevance-adequacy}), not arbitrary-premise local
consequence on every broad ternary frame.  Variable sharing belongs to
constant-free theorem implications in E and R
(Theorem~\ref{thm:p08-relevance-variable-sharing}); full RM has the explicit
exception and weaker replacement in
Theorem~\ref{thm:p08-rm-variable-sharing}.  Ackermann's gamma is a cited
admissible \emph{theoremhood metarule} for E and R
(Theorem~\ref{thm:p08-gamma-admissibility}), not an extra primitive sheet
rule from open assumptions \cite{AndersonBelnap1975,MeyerDunn1969,
DunnRestall2002}.  Full-language undecidability, the decidable pure
implicational fragments, and full RM decidability retain their distinct
targets in Theorems~\ref{thm:p08-relevance-undecidable},
\ref{thm:p08-relevance-implicational-decidable}, and
\ref{thm:p08-rm-decidable}.

Da Costa's \(C_1\), Sheet~\ref{sheet:hilbert-da-costa-c1}, is a different
Hilbert presentation with its exact postulates (1)--(15) and MP.
Its formula-level well-behavedness abbreviation is
Definition~\ref{def:p08-da-costa-c1}; it is neither a matrix value nor an
assertion that an entire theory is consistent.  The unrestricted and
controlled-explosion claims are the cited
Theorem~\ref{thm:p08-da-costa-c1-nonexplosion}, \cite{daCosta1974}.
LP instead obtains nonexplosion through designated-value preservation in
its fixed matrix.  The comparison in Table~\ref{tab:p08-paraconsistent-systems}
does not equip LP with a new consistency connective or identify the two
consequence relations.

The operational default-extension sheet,
\ref{sheet:reiter-default-extension}, is also distinct from a monotonic
Hilbert closure.  Prerequisites are tested in the current grounded stage;
justification consistency is tested against the fixed candidate extension;
the candidate is an extension only if it equals the resulting least grounded
closure.  Skeptical and credulous consequence quantify differently over
extensions.  Definition~\ref{def:p08-default-theory} and
Entry~\ref{entry:p08-reiter-default} supply the controlling construction
\cite{Reiter1980}.  The sheet is an operational fixed-point specification,
not an unconditional inference from a default to its conclusion.

\showcalculus{relevance-e}
\showcalculus{relevance-r}
\showcalculus{relevance-rm}
\showcalculus{hilbert-da-costa-c1}
\showcalculus{reiter-default-extension}

\section{Linear, Lambek, and other substructural calculi}
\label{sec:p14-nonclassical-substructural}

A resource context records how assumptions may be combined, reused,
discarded, or reordered.  The context shape is part of the calculus, not a
printing preference.  Definition~\ref{def:p08-context-consequence-properties}
and Table~\ref{tab:p08-consequence-properties} distinguish context-sensitive
cut, weakening and contraction from properties of a set-based consequence
relation.  The three relevant Hilbert systems above are not silently
replaced by resource sequents.

\begingroup\small
\begin{longtable}{@{}p{.18\linewidth}p{.36\linewidth}p{.38\linewidth}@{}}
\caption{Resource sheets: context policy and the exact metatheorem boundary.}
\label{tab:p14-nonclassical-resources}\\
\toprule
Sheet & Contexts, connectives, and side conditions & Cut, adequacy, and separation from variants \\
\midrule\endfirsthead
\caption[]{Resource sheets: context policy and the exact metatheorem boundary (continued)}\\
\toprule Sheet & Contexts, connectives, and side conditions & Cut, adequacy, and separation from variants \\
\midrule\endhead
\midrule\multicolumn{3}{r}{\emph{continued on next page}}\\
\endfoot
\endlastfoot
IMLL, \ref{sheet:imll} & Finite multisets, possibly empty, and exactly one
succedent. Tensor, unit and linear implication; exchange is built into
multiset equality, weakening and contraction are absent. Tensor-R splits
contexts rather than copying the same context into both premises. &
Identity and Cut are explicit. The source adequacy/cut statement remains
sketched in Entry~\ref{entry:p08-imll}, Annex~\ref{annex:p8.imll.source},
for the multiplicative core and its commutative residuated partially ordered
monoids. It does not cover exponentials. \\
\addlinespace
Source bang extension, \ref{sheet:ill-exponential-extension} & Same
single-succedent multiset calculus, extended by exactly four bang rules.
Weakening and contraction apply only to bang-prefixed antecedent formulas;
promotion requires the entire antecedent bang-prefixed. & Cut is inherited
as a rule. This exact source-bounded extension has no supplied adequacy,
cut-elimination or decision proof; it is not relabelled full ILL, and no
core or full-CLL metatheorem is transferred to it. \\
\addlinespace
One-sided CLL, \ref{sheet:cll-one-sided} & Finite one-sided multisets,
involutive De Morgan linear negation, multiplicatives, additives, units,
bang and quest. Additive with uses the same context in both premises;
tensor splits it. Structural exceptions are quest-prefixed, and promotion
requires an entirely quest-prefixed context. There is no introduction row
for additive zero. & The cited cut theorem is
Theorem~\ref{thm:p08-cll-cut-elimination}; full phase adequacy is
Theorem~\ref{thm:p08-phase-adequacy}, with the exponential phase conditions
of Definition~\ref{def:p08-phase-space}, \cite{Girard1995}.
The full, MALL, and multiplicative decision problems remain separate in
Theorem~\ref{thm:p08-linear-decision-boundaries}. \\
\addlinespace
Pfenning's full ILL surface, \ref{sheet:p14-ill-pfenning} & Two contexts:
unrestricted assumptions in \(\Gamma\), linear occurrences in \(\Delta\),
and one conclusion. Multiplicatives, additives, both units, exponentials,
unrestricted implication and first-order quantifiers are all displayed;
the linear contexts split only in the multiplicative two-premise rules. &
The exact cut-free sequent surface of the 27 January 1998 handout,
\cite[\S2.6]{Pfenning1998ILLSequent}. Its normal-deduction comparison is
source-stated, not an imported full-CLL metatheorem. The distinct inherited
four-rule bang extension is retained above without relabelling. \\
\addlinespace
Commutative affine ILL, \ref{sheet:affine-ill} & The IMLL grammar and
single-succedent multiset shape, with unrestricted antecedent weakening
added and contraction still absent. & Cut is inherited as a rule.
Definition~\ref{def:p08-affine-family-boundary} asserts no independent
adequacy or cut-elimination theorem for this selected extension,
\cite{Lincoln1995}. \\
\addlinespace
Associative Lambek L, \ref{sheet:lambek-l} & Finite \emph{nonempty ordered}
antecedents and one succedent; product and left/right divisions, without
unit, exchange, weakening or contraction. Context fragments may be empty
only while every displayed total antecedent is nonempty. Division-R
retains a nonempty residual context. &
Theorem~\ref{thm:p08-lambek-adequacy} cites residuated-semigroup completeness
and cut elimination, \cite{Buszkowski2010}. The nonempty restriction is
Lambek's L, not empty-antecedent L* or unit-bearing L1. The exact decision
classification is Entry~\ref{entry:p08-lambek-l}, \cite{Pentus2006}. \\
\bottomrule
\end{longtable}\endgroup

The algebraic difference matches the context difference.  IMLL uses a
commutative residuated partially ordered monoid, with a unit interpreting
the empty multiset.  Lambek L uses a residuated semigroup without a required
unit, and the left and right residuals may differ.  Their defining clauses
are Definition~\ref{def:p08-lambek-semantics} and
Definition~\ref{def:p10-residuated-lattice}.  Replacing an ordered context by
a multiset adds exchange; erasing the nonempty-antecedent condition changes
the allowed division introductions.  Neither change follows merely from a
shared symbol for product \cite{Lambek1958,GalatosJipsenKowalskiOno2007}.

In CLL, the displayed \(\bot\) is a multiplicative unit and the displayed
\(\top\) an additive unit; their rows are not classical explosion rules.
The interpretation of \(A\multimap B\) by linear negation and par belongs
to that one-sided language.  It does not identify its proof objects with the
single-succedent IMLL sheet.  Focusing is the separate polarity treatment in
Definition~\ref{def:p08-focusing-polarity}, \cite{Andreoli1992}; the presence
of a phase-completeness theorem alone does not specify a focused proof-search
procedure.

\showcalculus{imll}
\showcalculus{ill-exponential-extension}
\showcalculus{cll-one-sided}
\showcalculus{p14-ill-pfenning}
\showcalculus{affine-ill}
\showcalculus{lambek-l}

\section{Diagrammatic, algebraic, categorical, and program-logic rule collections}
\label{sec:p14-nonclassical-specialized}

\subsection{Equations, categorical diagrams, and contextual judgments}

An equational deduction, a commuting diagram, and a contextual predicate
sequent assert different things.  Sheet~\ref{sheet:equational-logic} treats
premises as \emph{universally interpreted identities} in nonempty algebras
over a fixed finitary signature.  Its substitution rule is sound for that
global identity interpretation; it is not licensed by a premise equation
holding under one chosen valuation.  The distinction is worked in
Example~\ref{def:p10-identities-valuations}.  The term algebra must be
nonempty, so an empty variable set with no nullary symbols is not silently
admitted.  The complete quotient-model proof and equivalence of compatibility
with contextual term replacement are
Theorem~\ref{thm:p10-equational-completeness},
Annex~\ref{annex:p10.eq.2}, reconstructed against
\cite[II\S14]{BurrisSankappanavar2012}.

\begingroup\small
\begin{longtable}{@{}p{.20\linewidth}p{.35\linewidth}p{.37\linewidth}@{}}
\caption{Algebraic and categorical sheet interfaces, without erasing contexts or types.}
\label{tab:p14-nonclassical-categorical}\\
\toprule
Registered surface & Judgment and essential typing & Exact result and proof route \\
\midrule\endfirsthead
\caption[]{Algebraic and categorical sheet interfaces, without erasing contexts or types (continued)}\\
\toprule Registered surface & Judgment and essential typing & Exact result and proof route \\
\midrule\endhead
\midrule\multicolumn{3}{r}{\emph{continued on next page}}\\
\endfoot
\endlastfoot
Equational logic, \ref{sheet:equational-logic} & Finite deductions from
universally interpreted identities; reflexivity, symmetry, transitivity,
compatibility and simultaneous term substitution. Nullary compatibility has
no premises. & Theorem~\ref{thm:p10-equational-completeness},
Annex~\ref{annex:p10.eq.2}, proves arbitrary-set equational consequence and
its finite-support property. This is completeness, not termination or
confluence of an oriented rewrite system. \\
\addlinespace
Pitts contextual first order, \ref{sheet:contextual-first-order} & External
sorted term equations and predicate sequents have separate judgment types.
The variable list \(\Gamma\) and finite hypothesis list \(\Phi\) are
separate; unused variables are retained. Substitution is simultaneous and
capture-avoiding; witness terms have their declared sorts, and quantifier
eigenvariables are fresh for the conclusion's retained hypotheses. &
Definition~\ref{def:p10-contextual-first-order} fixes the exact Figures
1--3, 5, 7 presentation. Soundness in the stated hyperdoctrine is
Theorem~\ref{thm:p10-doctrine-soundness},
Annex~\ref{annex:p10.quant-fol.2}. Sort objects may be uninhabited;
no context-erasing completeness or cut-elimination theorem is asserted
there, \cite{Pitts1995Draft}. \\
\addlinespace
Pitts adjoint presentation, \ref{sheet:contextual-adjoint-first-order} &
Fifteen directed inferences from the seven bi-rules, used over the same
equational/structural/theory-axiom basis. Upward directions with two upper
sequents are distinct; false has only its substantive downward direction.
Retained hypotheses and variable parameters remain exactly displayed. &
Theorem~\ref{thm:p10-contextual-adjoint-equivalence},
Annex~\ref{annex:p10.adjoint.2}, proves interderivability of judgments.
Both presentations retain Cut; the theorem does not identify their proof
trees or prove cut elimination. \\
\addlinespace
Parameter-free equality, \ref{sheet:contextual-parameter-free-equality} &
The two Eq0 directions have exactly the displayed one- or two-variable
contexts, no suppressed parameter prefix, and an empty retained hypothesis
list. & Theorem~\ref{thm:p10-contextual-equality-reduction},
Annex~\ref{annex:p10.adjoint.3}, recovers parameterized equality over the
common basis \emph{with implication and universal quantification}.
Those additional rule hypotheses are essential to the stated reduction,
\cite[Remark~5.6.2]{Pitts1995Draft}. \\
\bottomrule
\end{longtable}\endgroup

The contextual Weaken row appends fresh \emph{variable declarations};
Exchange exchanges formula hypotheses.  They are not two interchangeable
ways to discard a context.  The predicates' syntactic transport theorem is
Theorem~\ref{thm:p10-contextual-predicate-transport},
Annex~\ref{annex:p10.adjoint.1}; it explicitly asserts no corresponding
general transport rule for external term equations.  Nor is there a rule
reflecting a proved equality formula into equality of term arrows.  The
empty-sort example following Theorem~\ref{thm:p10-doctrine-tarski} explains
why deleting an unused declaration can change validity.

Categorical diagrams in Section~\ref{sec:p10-categorical-logic} record typed
arrows, equal composites, and universal properties.  The STLC/CCC equality
comparison is Theorem~\ref{thm:p10-stlc-ccc-completeness},
Annex~\ref{annex:p10.lambda.3}; the internal language recovering a given
category is Theorem~\ref{thm:p10-internal-ccc-equivalence},
Annex~\ref{annex:p10.internal.2}.  These supply exact uses of typed
composition, substitution, and the named equational diagram.  A displayed
pullback square requires its universal factorization property, not merely
its drawn commutativity (Definition~\ref{def:p10-topos}).  None of these
diagrams is thereby a complete Euler/Venn, existential-graph, or string-
diagram inference calculus.  A Boolean decision diagram likewise represents
a Boolean function under a fixed variable order; its cited canonicity theorem
is Theorem~\ref{thm:p12-robdd-canonicity}, \cite{Bryant1986}, not an inference
rule for categorical diagrams.

\showcalculus{equational-logic}
\showcalculus{contextual-first-order}
\showcalculus{contextual-adjoint-first-order}
\showcalculus{contextual-parameter-free-equality}

\subsection{Program assertions, heap frames, and operational goals}

Program-logic sheets expose the data whose preservation is being proved.
A Hoare judgment has assertions and a command, a separation judgment also
has heap ownership and safety, and an SLD step has a selected goal atom,
a standardized-apart clause, and a computed substitution.  They are not
alternative spellings of a single formula consequence relation.

\begingroup\small
\begin{longtable}{@{}p{.17\linewidth}p{.39\linewidth}p{.36\linewidth}@{}}
\caption{Program and operational rule collections: what each side condition protects.}
\label{tab:p14-nonclassical-programs}\\
\toprule
Sheet & Judgment and nonoptional conditions & Correctness/completeness boundary \\
\midrule\endfirsthead
\caption[]{Program and operational rule collections: what each side condition protects (continued)}\\
\toprule Sheet & Judgment and nonoptional conditions & Correctness/completeness boundary \\
\midrule\endhead
\midrule\multicolumn{3}{r}{\emph{continued on next page}}\\
\endfoot
\endlastfoot
Hoare while, \ref{sheet:hoare-while} & The partial/total tag is retained
through consequence, sequencing and branching. Assertion implication means
validity in every integer state and every rigid-logical-parameter value.
Assignment uses capture-avoiding substitution. Total while requires a
nonnegative integer-valued variant and a fresh \emph{rigid} old-value
parameter; the body premise must already establish total correctness. &
Soundness of every displayed row is proved in
Theorem~\ref{thm:p09-hoare-soundness}, Annex~\ref{annex:p9.hoare.1}.
Cook's relative completeness is partial-only, with expressible required
assertions and valid implication premises
(Theorem~\ref{thm:p09-hoare-relative-completeness}), \cite{Cook1978}.
The decreasing-integer row is editorial and is not misattributed to
Apt's distinct Harel-indexed rule. \\
\addlinespace
Separation frame, \ref{sheet:separation-logic} & Partial-correctness heap
triples require safety as well as a terminating-result condition.
Separating conjunction splits a heap into disjoint owned parts. The command
must satisfy the precise locality/frame property, and no modified stack
variable may occur free in the framed assertion. &
Definition~\ref{def:p09-separation-logic} fixes locality;
Theorem~\ref{thm:p09-frame-rule} cites the frame theorem to
\cite{OHearnReynoldsYang2001}. The sheet is this frame-rule collection,
not a claim of a complete proof system for all heap assertions. \\
\addlinespace
SLD and ground SLDNF, \ref{sheet:sld-resolution} & Ordered goals have a
selected literal. The clause is standardized apart; its finite-term mgu,
with occurs check, is applied to the whole resolvent. A computed answer is
restricted to the initial goal variables. Negative success needs a complete
finite failed subsidiary tree and a ground atom; nonground selection
flounders, while divergence is not finite failure. &
Definitions~\ref{def:p09-sld} and~\ref{def:p09-sldnf} separate the
procedures. SLD adequacy is Theorem~\ref{thm:p09-sld-adequacy}; SLDNF's
completion/nonfloundering conditions are
Theorem~\ref{thm:p09-sldnf-completion}, both source-cited to
\cite{Lloyd1987}. \(\mathsf{not}\) is finite failure, not classical
negation; the registered language excludes logical equality atoms. \\
\bottomrule
\end{longtable}\endgroup

For a graph-shaped search object with a separate full local proof,
Sheet~\ref{sheet:tableaux-alc} provides the selected ALC completion
procedure.  It initializes the global TBox conditions at every root and
new witness, rechecks subset ancestor blocking as labels grow, disables
only existential generation at a blocked source, and continues Boolean
and universal propagation across existing edges.  Its success condition is
one complete open branch.  The exact editorial correctness and termination
proof is Theorem~\ref{thm:p09-alc-tableau},
Annex~\ref{annex:p9.alc.1}; its statement does not claim a transcription of
a body-unverified book calculus.  This graph procedure, the hybrid tableau,
and finite-matrix tableaux have different node contents, closure tests,
and witness policies, which remain on their respective sheets.

\showcalculus{hoare-while}
\showcalculus{separation-logic}
\showcalculus{sld-resolution}
\showcalculus{tableaux-alc}

\index{calculus sheets!nonclassical and specialized}
\index{theorem-only inference}\index{local consequence}
\index{calculus sheets!context policies}\index{calculus sheets!proof boundaries}
